\documentclass{beamer}

\usetheme{Berlin}
\usecolortheme{Orchid}
\useoutertheme{miniframes}
\setbeamertemplate{navigation symbols}{}

\usepackage[utf8]{inputenc}
\usepackage[T1]{fontenc}
\usepackage{amsmath}
\usepackage{amssymb}
\usepackage{booktabs}
\usepackage{graphicx}
\graphicspath{{../images/}}

\usepackage{xcolor}
\usepackage{listings}
\usepackage{hyperref}
\usepackage{tikz}
\usetikzlibrary{arrows.meta,positioning,shapes.geometric,calc}

\lstset{
  language=Java,
  basicstyle=\ttfamily\small,
  keywordstyle=\color{blue},
  commentstyle=\color{gray},
  stringstyle=\color{teal},
  showstringspaces=false,
  columns=fullflexible,
  breaklines=true
}

% Per-slide source line (references shown on the slide itself).
\newcommand{\slidesources}[1]{%
  \vfill
  {\tiny\color{gray}\raggedright\sloppy\parbox{\linewidth}{Sources: #1}\par}%
}

% Unit-wise source strings for this subject (used by slides/notes).
% Path is relative to the lecture `latex/` folder (the compilation CWD).
% OOPS with Java (CSEG1044) — unit-wise sources
%
% These are short, student-facing reference strings intended to be used with:
%   - \slidesources{...}  (in beamer slides)
%   - \notesources{...}   (in lecture notes)
%
% Style inspiration: other subjects in My-Subjects/ use semicolon-separated sources
% and include an "accessed" date for web documentation.

\newcommand{\OOPSUnitISources}{Schildt, \textit{Java: A Beginner's Guide} (10e), McGraw-Hill (Java fundamentals + OOP basics).; Oracle Java Tutorials: Getting Started + Language Basics + Classes and Objects (accessed \today).; Java SE API docs (e.g., \texttt{String}, \texttt{Scanner}) (accessed \today).}

\newcommand{\OOPSUnitIISources}{Schildt, \textit{Java: The Complete Reference} (12e), McGraw-Hill (classes, inheritance, packages, interfaces).; Bloch, \textit{Effective Java} (3e), Addison-Wesley, 2018 (design choices; interfaces vs abstract classes).; Oracle Java Tutorials: Classes and Objects + Interfaces + Packages (accessed \today).; Java SE API docs (e.g., \texttt{Comparable}, \texttt{Runnable}) (accessed \today).}

\newcommand{\OOPSUnitIIISources}{Schildt, \textit{Java: The Complete Reference} (12e) (nested/inner classes, exceptions, threads, I/O).; Oracle Java Tutorials: Nested Classes + Exceptions + Concurrency + Basic I/O (accessed \today).; Java SE API docs (e.g., \texttt{Thread}, \texttt{AutoCloseable}) (accessed \today).}

\newcommand{\OOPSUnitIVSources}{Schildt, \textit{Java: The Complete Reference} (12e) (generics, lambdas, Swing, JDBC).; Bloch, \textit{Effective Java} (3e) (generics + lambdas).; Oracle Java Tutorials: Generics + Lambda Expressions + Swing + JDBC Basics (accessed \today).; Java SE API docs (e.g., \texttt{java.util.function}, \texttt{javax.swing}, \texttt{java.sql}) (accessed \today).}

\newcommand{\OOPSUnitVSources}{Schildt, \textit{Java: The Complete Reference} (12e) (Collections Framework + wrapper classes).; Oracle Java docs: Collections Framework overview (accessed \today).; Java SE API docs (\texttt{java.util} collections; wrapper classes in \texttt{java.lang}) (accessed \today).}

\newcommand{\OOPSUnitVISources}{Git SCM: \textit{Pro Git} (2e) (version control workflow), accessed \today.; GitHub Docs (repositories, issues, pull requests), accessed \today.; Oracle Java Tutorials: Swing + JDBC (accessed \today).; JUnit 5 User Guide (accessed \today).; Apache Maven Project: Getting Started (accessed \today).; Gradle User Manual: Getting Started (accessed \today).; UML class diagrams: UML 2.x overview/spec (accessed \today).}

\newcommand{\OOPSMiscWhyInterfacesSources}{Schildt, \textit{Java: The Complete Reference} (12e) (interfaces).; Bloch, \textit{Effective Java} (3e) (prefer interfaces where appropriate).; Oracle Java Tutorials: Interfaces (accessed \today).; Java SE API docs (e.g., \texttt{Comparable}, \texttt{Runnable}) (accessed \today).}



\title[OOPS with Java]{Object Oriented Programming (CSEG1044)}
\subtitle{Unit 01 -- Lecture 03: I/O + Command Line + Types + Operators}
\begin{document}

\begin{frame}[fragile]
  \titlepage
  \vspace{0.5em}
  \slidesources{\OOPSUnitISources}
\end{frame}

\begin{frame}{Quick Links}
  \centering
  \hyperlink{sec:overview}{\beamerbutton{Overview}}\hspace{0.6em}
  \hyperlink{sec:concepts}{\beamerbutton{Key Topics}}\hspace{0.6em}
  \hyperlink{sec:demo}{\beamerbutton{Demo}}\hspace{0.6em}
  \hyperlink{sec:summary}{\beamerbutton{Summary}}
  \slidesources{\OOPSUnitISources}
\end{frame}

\begin{frame}{Agenda}
  \tableofcontents
  \slidesources{\OOPSUnitISources}
\end{frame}

\section{Overview}
\label{sec:overview}

\begin{frame}{Lecture Snapshot}
\begin{itemize}[<+->]
  \item \textbf{Unit focus:} Introduction to OOPs and Java Fundamentals
  \item \textbf{Course thread:} Use I/O + Command Line + Types + Operators to make the first text-based version of Campus Resource Manager easier to reason about before the full object model arrives.
\end{itemize}
  \slidesources{\OOPSUnitISources}
\end{frame}

\begin{frame}{Recall Before You Start}
\begin{itemize}[<+->]
  \item \textbf{Program structure}: You already know where execution starts in Java; this lecture fills that structure with input and expressions. Main reference: \texttt{Unit 01 / Lecture 02 slides/notes}.
  \item \textbf{Basic data handling}: The C prerequisite gives the intuition for values and operators; Java adds stronger typing rules. Main reference: Programming in C prerequisite.
\end{itemize}
  \slidesources{\OOPSUnitISources}
\end{frame}


\begin{frame}{Learning Outcomes}
  \begin{itemize}[<+->]
    \item Read input using \texttt{Scanner} and print formatted output.
    \item Use command line arguments (\texttt{args}) for simple programs.
    \item Apply primitive types, operators, precedence, and casting correctly.
  \end{itemize}
  \slidesources{\OOPSUnitISources}
\end{frame}

\begin{frame}{Key Terms to Know}
\begin{itemize}[<+->]
  \item \textbf{Token vs line}: token = one word/number; line = full line of text.
  \item \textbf{Formatting}: printing numbers/text with a chosen format (decimal places, alignment).
  \item \textbf{Primitive types}: built-in value types like \texttt{int}, \texttt{double}, \texttt{char}, \texttt{boolean}.
  \item \textbf{Type casting}: converting a value from one type to another (widening/narrowing).
  \item \textbf{Precedence}: order in which operators are evaluated.
\end{itemize}
  \slidesources{\OOPSUnitISources}
\end{frame}

\begin{frame}{Study Flow for This Lecture}
\begin{enumerate}[<+->]
  \item Refresh the prerequisite bridge before reading new ideas in isolation.
  \item Study the main build in this order: \textit{Console I/O: \texttt{Scanner} (most common for beginners)} $\rightarrow$ \textit{Console I/O: \texttt{BufferedReader} (fast, line-oriented)} $\rightarrow$ \textit{Output formatting} $\rightarrow$ \textit{Command line arguments}.
  \item Trace the worked example and connect each line back to one concept frame.
  \item Attempt the mini exercises before moving to the recap and exit check.
\end{enumerate}
  \slidesources{\OOPSUnitISources}
\end{frame}


\section{Key Topics}
\label{sec:concepts}

\begin{frame}{Unit 01: Introduction to OOPs}
  \begin{itemize}[<+->]
    \item Lecture focus: I/O + Command Line + Types + Operators
  \end{itemize}
  \slidesources{\OOPSUnitISources}
\end{frame}

\begin{frame}{Today: Roadmap}
  \begin{itemize}[<+->]
    \item Console I/O (`Scanner`, `BufferedReader`) and output formatting
    \item Command line arguments and use-cases
    \item Data types, variables, operators, precedence, type casting
  \end{itemize}
  \slidesources{\OOPSUnitISources}
\end{frame}

\begin{frame}{Console input: \texttt{Scanner} (beginner-friendly)}
  \begin{itemize}[<+->]
    \item Reads \textbf{tokens} from \texttt{System.in}: \texttt{nextInt()}, \texttt{nextDouble()}, \texttt{next()}.
    \item Reads \textbf{lines}: \texttt{nextLine()}.
    \item Common pitfall: mixing \texttt{nextInt()} then \texttt{nextLine()} (newline remains in buffer).
    \item Fix: consume leftover newline with one extra \texttt{nextLine()}.
  \end{itemize}
  \slidesources{\OOPSUnitISources}
\end{frame}

\begin{frame}{Output formatting: \texttt{printf}}
  \begin{itemize}[<+->]
    \item \texttt{println} is fine for simple output.
    \item \texttt{printf} is better for formatted numbers/text.
  \end{itemize}
  \begin{block}{Examples}
    \begin{itemize}
      \item \texttt{\%d} integer,\; \texttt{\%f} decimal,\; \texttt{\%.2f} 2 decimals,\; \texttt{\%s} string
      \item \texttt{\%n} portable newline
    \end{itemize}
  \end{block}
  \slidesources{\OOPSUnitISources}
\end{frame}

\begin{frame}{Command line arguments (\texttt{args})}
  \begin{itemize}[<+->]
    \item Run: \texttt{java Sum 10 20 30}
    \item Then: \texttt{args[0] = "10"}, \texttt{args[1] = "20"}, ...
    \item Parse numbers: \texttt{Integer.parseInt(args[i])}
    \item Pitfall: invalid input throws \texttt{NumberFormatException}
  \end{itemize}
  \slidesources{\OOPSUnitISources}
\end{frame}

\begin{frame}{Primitive types (quick)}
  \begin{itemize}[<+->]
    \item Integers: \texttt{byte, short, int, long}
    \item Floating-point: \texttt{float, double}
    \item Others: \texttt{char} (Unicode), \texttt{boolean}
  \end{itemize}
  \slidesources{\OOPSUnitISources}
\end{frame}

\begin{frame}{Operator pitfalls (must know)}
  \begin{itemize}[<+->]
    \item Integer division: \texttt{5/2 == 2} (use \texttt{5/2.0})
    \item Precedence: multiplication before addition (use parentheses when unsure)
    \item Casting: narrowing needs explicit cast (may lose data)
    \item Short-circuit: \texttt{\&\&} and \texttt{||} help safe null checks
  \end{itemize}
  \slidesources{\OOPSUnitISources}
\end{frame}

\begin{frame}{Checkpoint and Reflection}
\begin{block}{Reflection prompt}
Where would \textit{I/O + Command Line + Types + Operators} matter most in Campus Resource Manager, and which design choice would you justify using this lecture's ideas?
\end{block}
\begin{block}{Quick self-check}
Which part needs one more example before you move on: \textit{Console I/O: \texttt{Scanner} (most common for beginners)}, \textit{Console I/O: \texttt{BufferedReader} (fast, line-oriented)}, the worked example, or the recap?
\end{block}
  \slidesources{\OOPSUnitISources}
\end{frame}

\section{Demo / Example}
\label{sec:demo}

\begin{frame}[fragile]{Example: Args + Scanner + Formatting (1/2)}
\begin{lstlisting}
import java.util.Scanner;

public class InputDemo {
    public static void main(String[] args) {
        int maxMarks = 100;
        if (args.length > 0) maxMarks = Integer.parseInt(args[0]);

        Scanner sc = new Scanner(System.in);
        System.out.print("Enter your name: ");
        String name = sc.nextLine();
\end{lstlisting}
  \slidesources{\OOPSUnitISources}
\end{frame}

\begin{frame}[fragile]{Example: Args + Scanner + Formatting (2/2)}
\begin{lstlisting}
// continued...
        System.out.print("Enter marks (0-" + maxMarks + "): ");
        int marks = sc.nextInt();

        double percent = marks * 100.0 / maxMarks;
        System.out.printf("Hello %s, marks=%d/%d, percent=%.2f%%%n",
                name, marks, maxMarks, percent);
        sc.close();
    }
}
\end{lstlisting}
  \slidesources{\OOPSUnitISources}
\end{frame}

\begin{frame}{Mini Exercises (2--3 mins)}
  \begin{enumerate}[<+->]
    \item Predict output: \texttt{System.out.println(5/2);}
    \item How to get 2.5 from 5 and 2?
    \item If you run \texttt{java Demo A B}, what is \texttt{args.length}?
  \end{enumerate}
  \slidesources{\OOPSUnitISources}
\end{frame}

\begin{frame}{Watch-outs}
\begin{itemize}[<+->]
  \item Mixing numeric and line-based input without understanding Scanner behavior.
  \item Ignoring operator precedence and misreading the result.
  \item Expecting Java to convert unrelated types automatically.
\end{itemize}
  \slidesources{\OOPSUnitISources}
\end{frame}

\section{Summary}
\label{sec:summary}

\begin{frame}{Key Takeaways}
  \begin{itemize}[<+->]
    \item \texttt{Scanner} is easy for input; know the \texttt{nextInt()} vs \texttt{nextLine()} pitfall.
    \item \texttt{printf} makes output readable (format specifiers).
    \item Command line args are strings; parse to numbers when needed.
    \item Watch integer division, precedence, and casting.
  \end{itemize}
  \slidesources{\OOPSUnitISources}
\end{frame}

\begin{frame}{Checkpoint Questions}
  \begin{enumerate}[<+->]
    \item Why can \texttt{nextLine()} return empty after \texttt{nextInt()}?
    \item What is the type of \texttt{args[i]}?
    \item How do you avoid integer division when you want decimals?
  \end{enumerate}
  \slidesources{\OOPSUnitISources}
\end{frame}

\begin{frame}{Next Study Steps}
\begin{itemize}[<+->]
  \item Revisit the recall references if any older idea still feels weak.
  \item Rework the lecture example without looking at the notes first.
  \item Attempt the self-study practice and then answer the exit questions in full sentences.
  \item Apply the idea once inside Campus Resource Manager to make the concept stick.
\end{itemize}
  \slidesources{\OOPSUnitISources}
\end{frame}

\begin{frame}{Next Lecture Preview}
  \textbf{Next:} Control Flow + Arrays + Strings
  \vspace{0.6em}
  \begin{itemize}[<+->]
    \item Bring one example where you used a loop in C.
    \item We will solve simple array + string problems together.
  \end{itemize}
  \slidesources{\OOPSUnitISources}
\end{frame}

\end{document}
